\documentclass[onecolumn, hidelinks, 12pt]{article}
\usepackage[hmargin={0.5in, 0.5in}, vmargin={0.8in, 0.8in}]{geometry}
\usepackage{booktabs}
%\usepackage{hyperref}
%\usepackage[anythingbreaks]{breakurl}
\usepackage[hyphens,spaces,obeyspaces]{url}
\usepackage{hyperref}

\newcommand{\pkg}[1]{{\normalfont\fontseries{b}\selectfont #1}}
\let\proglang=\textsf
\let\code=\texttt


\usepackage{enumitem}
\setlist{parsep=0pt, leftmargin=4mm, topsep=0pt, itemsep=8pt}
% \setlist{nolistsep}

\usepackage{fancyhdr}
\pagestyle{fancy}
\renewcommand{\headrulewidth}{0.4pt}
\renewcommand{\footrulewidth}{0.4pt}

\lhead{\sf CCLA: Data Science}
\rhead{Fall 2023 --- Spring2024}
\lfoot{}
\rfoot{} 


\begin{document}

\title{CCLA: Data Science, 2023--2024}
\date{} %
\maketitle

\thispagestyle{fancy}


\begin{description}
\item[Instructors:] \hspace{0pt}
  
  HaiYing Wang, Associate Professor, Department of Statistics, University of Connecticut\\
  Jun Yan, Professor, Department of Statistics, University of Connecticut\\
  
\item[Lectures:] \hspace{0pt}

  Location: TBA
  
  Sunday: 11:00--11:50 ET
  

\item[Course Description:]
Data science is a fast developing science of extracting meaningful
information from massive data for better decision making. It is
interdisciplinary by nature, involving statistics, computing, and
domain knowledge. Important principles of data science will be
elaborated through interactive simulations, games and examples.

\item[Prerequisites:]
  Some exposure to programming and data science would be a plus,
  but is not required.  

\item[Topics:] \hspace{0pt}
  
\begin{enumerate} % [noitemsep]
\item
  Introduction to Julia, random number generation, data summaries
\item
  Exploratory data analysis, visualization
\item
  Paradoxes explained
\item
  At the infinity horizon (as sample size increases)
\item
  Hide-and-seek: Estimating the unknowns
\item
  Sampling design, data collection
\item
  Correlation and regression
\item
  Hypothesis testing
\item
  A real-world data science project
\end{enumerate}

\item[Computing:]
  Please have Julia preinstalled:
  \url{https://julialang.org}

\item[References:] \hspace{0pt}
  
  % Some are open-source books while others are not.
  \begin{itemize}
  \item
    Joshi, A. (2016). Julia for Data Science. Packt Publishing Ltd.
  \item
    Storopoli, J., Huijzer, R., \& Alonso, L. (2022). Julia Data Science. Editeur indépendant.
    \url{https://github.com/JuliaDataScience/JuliaDataScience}
  \item
    Nazarathy, Y., \& Klok, H. (2021). Statistics with Julia: Fundamentals for data science, machine learning and artificial intelligence. Springer Nature.
    \url{https://github.com/h-Klok/StatsWithJuliaBook}
  \end{itemize}
  

\end{description}

\end{document}
